\chapter{Risultati e Caso di Studio}
\label{chap:risultati}

% Questo capitolo è la vetrina del tuo lavoro. Usa molte immagini (screenshot).

\section{Il Sito Web Risultante}
% Mostra degli screenshot delle pagine principali del sito in modalità di visualizzazione normale (non-editing).
% Es. la home page, la pagina `people`, la pagina `publications`.

\section{L'Esperienza di Editing}
% Mostra degli screenshot del sito in modalità modifica (`#simply-edit`).
% - Evidenzia le aree editabili (con il contorno tratteggiato di SimplyEdit).
% - Mostra la toolbar principale in alto.
% - Mostra le toolbar contestuali che appaiono quando selezioni un testo o un'immagine.
% - Mostra la sitemap di SimplyEdit per la creazione di nuove pagine.

\section{Il Workflow di Versioning in Azione}
% Questa è la prova che il tuo sistema funziona come progettato.
% Fai uno screenshot del terminale che esegue `git log`.
% Evidenzia i messaggi di commit generati automaticamente dal tuo server Node.js (es. "Aggiornamento contenuti sito via SimplyEdit", "Aggiunta nuova pagina: /people.html").
% Questo dimostra visivamente il successo dell'integrazione con Git.

\section{Ricerca di Pubblicazioni Esterne}
% Mostra uno screenshot della pagina delle pubblicazioni dopo aver effettuato una ricerca.
% Fai vedere la tabella dei risultati popolata con i dati recuperati da `art.torvergata.it`.
% Questo dimostra il funzionamento dell'integrazione con servizi esterni.

\section{Sistema di Autenticazione}
% Mostra uno screenshot della pagina di login (`login.html`).
% Spiega brevemente come l'accesso alla modalità di modifica sia precluso senza aver prima effettuato il login.
